\section{Introduction}
\label{sec:introduction}

Binary optimization asks for an assignment of finitely many decisions that
minimizes a cost or maximizes a reward. The variables may describe the two
sides of a graph cut, selected facilities, accepted constraints, or a spin
configuration. A common mathematical representation is quadratic
unconstrained binary optimization (QUBO); after an affine change of
variables, the same objective is an Ising energy. The representation is
general enough to encode many problems, but it does not remove their
combinatorial difficulty: $N$ unrestricted binary variables have $2^N$
possible assignments. MaxCut and general Ising optimization already contain
NP-hard cases \citep{karp1972,barahona1982,lucas2014}.

For MaxCut, the input is an undirected weighted graph $G=(V,E,w)$ and a
feasible solution is a partition $S\subseteq V$ of its vertices. The problem
is to find
\begin{equation}
 \max_{S\subseteq V}\; C(S),\qquad
 C(S)=\sum_{\{i,j\}\in E}w_{ij}
 \mathbf 1\{ |\{i,j\}\cap S|=1\}.
 \label{eq:maxcut-statement}
\end{equation}
Equivalently, a binary spin $s_i\in\{-1,+1\}$ marks each vertex's side of
the partition. In this paper, a reported MaxCut solution is such a binary
partition, scored with the original edge weights; the continuous QeFEM
state is the means of finding it.

Many useful solvers therefore exchange exhaustive enumeration for a search
over an easier auxiliary space. Simulated annealing explores binary states
stochastically \citep{kirkpatrick1983}. The free-energy machine (\FEM)
optimizes a factorized thermal distribution with gradients
\citep{shen2025}. Local quantum annealing (\LQA) instead optimizes a pure
product-state surrogate for a changing transverse-field Hamiltonian
\citep{bowles2022}. Discrete simulated bifurcation (\dSB) updates classical
position and momentum variables with a sign-dependent interaction
\citep{goto2021dsbm}. These algorithms differ in objective, dynamics, and
computational budget. Their common feature is that a continuous or
population-based search eventually returns an ordinary binary assignment.

\QeFEM{} places the search in a product of single-qubit mixed states. A
local state has a longitudinal component, which records the expected binary
decision; a transverse component, which couples to a driver; and an entropy,
which resists premature polarization. The objective is the mean-field free
energy of a transverse-field Ising model. We minimize it with classical
tensor operations and independently initialized replicas. The word
``quantum'' identifies the state family and the noncommuting driver in the
surrogate Hamiltonian. It does not imply entanglement, physical quantum
evolution, or quantum computational speedup.

The point of a detailed derivation is to identify precisely what this
relaxation can and cannot promise. The exact Gibbs state is generally
intractable, so the product ansatz introduces a substantial approximation.
The resulting free energy is nonconvex at low temperature. Annealing and
replica search can find strong states, but neither follows an adiabatic
theorem, and neither guarantees the optimal cut. Furthermore, the later
G-set experiments use a deliberately modified, discrete-neighbour search
direction. That update shares the state representation and final score with
the smooth method, but it is not the exact gradient of the original free
energy. Treating the two variants as identical would obscure both the
mathematics and the evidence.

The results require similar care. The main saved comparison consists of
frozen \QeFEM{} and \LQA{} configurations on K2000, Wishart, and Chook
instances. The G-set result combines historical attainment with a subsequent
adaptive search and fresh-seed confirmation. A separate FEM notebook
contains informative K2000 runs of \FEM{} and \dSB, but those runs are not
paired with the saved \QeFEM{} experiment in compute, numerical platform,
or seed protocol. We preserve these distinctions throughout rather than
compressing unlike measurements into one apparent league table.

This expanded manuscript is organized around four questions. First, how
does a binary Ising objective lead to the \QeFEM{} free energy? Second, what
do its stationary equations, gradients, and instability scale say about the
search dynamics? Third, what was actually run, including optimizer settings,
replica counts, seeds, and readout? Fourth, what do the stored endpoints
establish, and where does the method still fall short? The figures are
chosen to answer these questions directly: a winner trajectory and entropy
trace expose the annealing dynamics; population scaling diagnoses the K2000
plateau; instance-level curves show where planted error changes; and a
gap-before/after plot displays which revisited G-set graphs improved and
which supplied references were newly attained.
